\documentclass[12pt,a4paper]{extreport}

\usepackage{fontspec}
\usepackage[english,russian,armenian]{babel}

% Шрифт
\setmainfont{DejaVu Serif}

% 1. Поля страницы
\usepackage[a4paper, left=3.0cm, right=1.5cm, top=2.0cm, bottom=2.0cm]{geometry}

% 2. Межстрочный интервал (1.5)
\usepackage{setspace}
\onehalfspacing

% 3. Абзацный отступ и переносы
\usepackage{microtype}
\setlength{\parindent}{1.25cm}
\emergencystretch=3em
\tolerance=2000

% 4. Пакеты для содержания (формулы, рисунки, таблицы)
\usepackage{amsmath}
\usepackage{graphicx}
\usepackage{amssymb}

% 5. Заголовки (titlesec)
\usepackage{titlesec}

% Глава: «1. Название» по центру, жирным
\titleformat{\chapter}[hang]
  {\normalfont\Large\bfseries\centering}
  {\thechapter.}
  {0.5em}
  {}
\titlespacing*{\chapter}{0pt}{0pt}{20pt}

% Раздел: «1.1. Название» слева
\titleformat{\section}[hang]
  {\normalfont\large\bfseries}
  {\thesection.}
  {0.5em}
  {}
\titlespacing*{\section}{0pt}{12pt}{6pt}

% Подраздел: «1.1.1. Название» слева
\titleformat{\subsection}[hang]
  {\normalfont\normalsize\bfseries}
  {\thesubsection.}
  {0.5em}
  {}
\titlespacing*{\subsection}{0pt}{10pt}{6pt}

\begin{document}

\chapter{Լուսացույցների կառավարման մեթոդների ակնարկ}

\section{Լուսացույցների կառավարման խնդիրը և որակի չափանիշները}

\subsection{Խնդրի բնութագիրը}

Խաչմերուկը քաղաքային ճանապարհային ցանցի այն հատվածն է, որտեղ հատվում են տարբեր ուղղություններով շարժվող տրանսպորտային հոսքերը։ Խաչմերուկի թողունակությունը սահմանափակ է, քանի որ կանաչ լույս միաժամանակ կարող է ստանալ միայն չհակասող ուղղությունների խումբը։ Լուսացույցի կառավարումը հանգում է հետևյալ հարցին՝ ո՞ր ուղղությանը և որքա՞ն ժամանակով տալ կանաչ լույս։ Սխալ որոշումները հանգեցնում են հերթերի առաջացման, ժամանակի կորուստների, վառելիքի ավելորդ ծախսի և արտանետումների աճի։ [աղբյուր]

Լուսացույցի ցիկլի երկարության ընտրությունը պայմանավորված է երկու հակադիր գործոններով։ Յուրաքանչյուր փուլի փոփոխության ժամանակ առաջանում է կորսված ժամանակ, որը կապված է դեղին լույսի և կանգնած մեքենաների արագացման հետ։ Ցիկլը կարճացնելիս այդ կորուստների մասնաբաժինը մեծանում է, և ժամանակի զգալի մասը չի օգտագործվում մեքենաների անցման համար։ Ցիկլը երկարացնելիս կորուստների մասնաբաժինը փոքրանում է, սակայն կարմիր լույսի ժամանակ կուտակվում են ավելի երկար հերթեր։ Այսպիսով, գոյություն ունի ցիկլի օպտիմալ երկարություն, որը կախված է հոսքի ինտենսիվությունից։ [աղբյուր՝ Վեբսթեր]

\subsection{Որակի չափանիշները}

Կառավարման որակը գնահատելու համար գրականության մեջ կիրառվում են մի քանի չափանիշներ՝

\begin{itemize}
  \item \textbf{սպասման ժամանակ (հապաղում)՝} ժամանակը, որի ընթացքում մեքենան կանգնած է եղել,
  \item \textbf{հերթի երկարություն} և կանգնած մեքենաների քանակ,
  \item \textbf{հոսքի միջին արագություն,}
  \item \textbf{թողունակություն՝} մեկ ժամվա ընթացքում խաչմերուկով անցած մեքենաների քանակ,
  \item \textbf{ճանապարհի ժամանակ։}
\end{itemize}

Այս չափանիշներից և ոչ մեկը առանձին չի բնութագրում կառավարման որակը ամբողջությամբ, քանի որ դրանք չափում են տարբեր երևույթներ։ Օրինակ, սպասման ժամանակը կարող է փոքր լինել այն դեպքում, երբ մեքենաների մի մասը չի հասցնում մտնել ցանց և չի հաշվառվում։ Մեր փորձերում բարձր բեռնվածության պայմաններում PPO-ն ցուցաբերում էր ավելի ցածր սպասման ժամանակ, քան actuated կառավարումը, սակայն ցանցի մուտքի մոտ մնում էր մոտ 400 մեքենա (actuated-ի դեպքում՝ մոտ 60)։ Այդ պատճառով աշխատանքում կիրառվում է չափանիշների համակարգ, որը ներառում է նաև ընդհանուր հերթը՝ կանգնած և ցանց մտնել սպասող մեքենաների գումարը։

\section{Լուսացույցների կառավարման դասական մեթոդները}

\subsection{Ֆիքսված ցիկլով կառավարումը և Վեբսթերի մեթոդը}

Ֆիքսված ցիկլով կառավարման դեպքում լուսացույցի փուլերի հաջորդականությունը և յուրաքանչյուր փուլի տևողությունը նախապես սահմանվում են և չեն փոխվում ժամանակի ընթացքում։ Ցիկլը բաղկացած է կանաչ, դեղին և (անհրաժեշտության դեպքում) բոլոր ուղղությունների համար կարմիր լույսի ժամանակահատվածներից։ Այս մոտեցման առավելությունը պարզությունն է. այն չի պահանջում դետեկտորներ և հեշտությամբ համակարգվում է հարևան խաչմերուկների հետ։ Թերությունը հոսքի փոփոխություններին չհարմարվելն է. պարամետրերը հաշվարկվում են միջին ինտենսիվության համար։ [աղբյուր]

Ցիկլի երկարության և փուլերի բաժանման դասական հաշվարկային մեթոդը առաջարկել է Ֆ. Վեբսթերը [աղբյուր՝ Վեբսթեր]։ Մեթոդում օգտագործվում են հետևյալ մեծությունները։ Յուրաքանչյուր փուլի համար որոշվում է հոսքի բեռնվածությունը
\begin{equation}
  y_i = \frac{q_i}{s_i},
\end{equation}
որտեղ $q_i$-ն փուլի կրիտիկական ուղղության իրական հոսքն է, իսկ $s_i$-ն՝ հագեցման հոսքը, այսինքն՝ այն առավելագույն քանակի մեքենաները ժամում, որոնք կարող են անցնել խաչմերուկով՝ շարունակական կանաչ լույսի դեպքում։ Խաչմերուկի ընդհանուր բեռնվածությունը
\begin{equation}
  Y = \sum_i y_i .
\end{equation}
Եթե $Y \ge 1$, խաչմերուկը գերբեռնված է ցիկլի ցանկացած երկարության դեպքում։ Վեբսթերը ցույց է տվել, որ միջին հապաղումը նվազագույն է ցիկլի մոտավորապես հետևյալ երկարության դեպքում.
\begin{equation}
  C_0 = \frac{1{,}5\,L + 5}{1 - Y},
\end{equation}
որտեղ $L$-ը ցիկլի ընթացքում կորսված ընդհանուր ժամանակն է։ Արդյունավետ կանաչ ժամանակը բաժանվում է փուլերի միջև ըստ բեռնվածության.
\begin{equation}
  g_i = \left(C_0 - L\right)\frac{y_i}{Y}.
\end{equation}
Այսպիսով, ավելի բեռնված ուղղությունը ստանում է ավելի երկար կանաչ լույս։ Մեթոդի սահմանափակումն այն է, որ այն օգտագործում է միջինացված հոսքեր և չի արձագանքում կարճաժամկետ տատանումներին։

\subsection{Actuated կառավարումը}

Actuated (դետեկտորներով կառավարվող) համակարգերում խաչմերուկի մոտեցումների վրա տեղադրվում են դետեկտորներ, և փուլի տևողությունը որոշվում է մեքենաների ներկայությամբ։ Յուրաքանչյուր փուլի համար սահմանվում են նվազագույն և առավելագույն կանաչ տևողություններ։ Նվազագույն ժամանակից հետո փուլը երկարացվում է այնքան, քանի դեռ դետեկտորը գրանցում է մեքենաների հոսք, և ընդհատվում է, երբ մեքենաների միջև բացը գերազանցում է սահմանված շեմը, կամ երբ հասնում է առավելագույն տևողությանը։ [աղբյուր]

Այս մոտեցումը պարզ է, չի պահանջում հոսքի մոդել և աշխատում է տեղային տվյալներով։ Նրա թերություններն են՝ խաչմերուկների միջև համակարգման բացակայությունը և պարամետրերի (նվազագույն, առավելագույն տևողություններ, բացի շեմ) ձեռքով ընտրությունը։

\subsection{Ադապտիվ համակարգեր՝ SCOOT և SCATS}

Քաղաքային մասշտաբով կիրառվող ադապտիվ համակարգերը օգտագործում են դետեկտորների տվյալները իրական ժամանակում և փոփոխում են կառավարման պարամետրերը ըստ իրավիճակի։ Ամենահայտնի համակարգերն են SCOOT-ը և SCATS-ը։

SCOOT (Split, Cycle and Offset Optimisation Technique) համակարգը մշակվել է Մեծ Բրիտանիայում 1980-ականների սկզբին [աղբյուր]։ Այն կառուցում է տրանսպորտային հոսքերի առցանց մոդել և փոքր քայլերով ճշգրտում է երեք պարամետր՝ փուլերի բաժանումը (split), ցիկլի երկարությունը (cycle) և հարևան խաչմերուկների լույսերի հերթափոխման շեղումը (offset)։

SCATS (Sydney Coordinated Adaptive Traffic System) համակարգը մշակվել է Սիդնեյում [աղբյուր]։ Այն ունի աստիճանային կառուցվածք. տեղային կարգավորիչները աշխատում են ըստ նախապես մշակված պլանների, իսկ կենտրոնական մակարդակը ընտրում է տվյալ պահին համապատասխան պլանը՝ ելնելով հագեցվածության աստիճանից։

Այս համակարգերը ցույց են տալիս, որ ադապտիվ կառավարումը գոյություն ունի վաղուց և կիրառվում է գործնականում։ Դրանց թերությունը կենտրոնացված կառուցվածքն է, ինչպես նաև կառավարման պարամետրերի կարգավորման բարդությունը։

\subsection{Max-Pressure մեթոդը}

Max-Pressure մեթոդը (առավելագույն ճնշման մեթոդ) առաջարկվել է Պ. Վարայայի կողմից [աղբյուր՝ Վարայա, 2013]։ Յուրաքանչյուր շարժման համար սահմանվում է ճնշումը՝ մուտքային գոտու հերթի և ելքային գոտու հերթի տարբերությունը։ Փուլի ճնշումը այդ փուլի շարժումների ճնշումների գումարն է.
\begin{equation}
  P(\varphi) = \sum_{(l,m)\in\varphi} \bigl(x_l - x_m\bigr),
\end{equation}
որտեղ $x_l$-ը մուտքային գոտու, իսկ $x_m$-ը ելքային գոտու հերթն է։ Յուրաքանչյուր պահի ընտրվում է առավելագույն ճնշմամբ փուլը։

Մեթոդի հիմնական առավելությունն այն է, որ նա ապակենտրոնացված է. որոշումը հիմնվում է միայն տեղային հերթերի վրա և չի պահանջում հոսքի ինտենսիվության գնահատում։ Վարայան ցույց է տվել, որ մեթոդը ապահովում է հերթերի կայունություն, երբ պահանջարկը գտնվում է խաչմերուկների ցանցի թողունակության շրջանում [աղբյուր՝ ստուգել ձևակերպումը]։ Մեկ մեկուսացված խաչմերուկի դեպքում, երբ ելքային գոտիները ազատ են, ճնշումը հանգում է մուտքային հերթերի համեմատությանը։

\subsection{Մեթոդների համեմատություն}

Աղյուսակ~\ref{tab:classic} ամփոփում է դիտարկված դասական մեթոդների հիմնական հատկությունները։

\begin{table}[h]
\centering
\small
\begin{tabular}{|p{3cm}|p{2.6cm}|p{2.6cm}|p{2.6cm}|p{2.6cm}|}
\hline
 & Ֆիքսված (Վեբսթեր) & Actuated & SCOOT / SCATS & Max-Pressure \\
\hline
Դետեկտորներ & ոչ & այո & այո & այո \\
\hline
Հոսքի մոդել & միջին հոսքեր & ոչ & այո (SCOOT) & ոչ \\
\hline
Կառուցվածք & տեղային & տեղային & կենտրոնացված & ապակենտրոնացված \\
\hline
Հարմարվողականություն & չկա & տեղային & ցանցային & տեղային \\
\hline
Տեսական երաշխիքներ & հաշվարկային & չկա & չկա & կայունություն \\
\hline
\end{tabular}
\caption{Դասական մեթոդների համեմատություն}
\label{tab:classic}
\end{table}

Դասական մեթոդները թափանցիկ են և կայուն, սակայն պահանջում են պարամետրերի ձեռքով կարգավորում և սահմանափակ ձևով են հաշվի առնում խաչմերուկների բարդ փոխազդեցությունը։ Այս սահմանափակումներն են պատճառը, որ վերջին տարիներին լուսացույցների կառավարման համար ուսումնասիրվում են ամրապնդող ուսուցման մեթոդները։

\section{Ամրապնդող ուսուցում}

Ամրապնդող ուսուցումը (reinforcement learning, RL) մեքենայական ուսուցման ուղղություն է, որտեղ գործակալը սովորում է որոշումներ ընդունել՝ փոխազդելով միջավայրի հետ և ստանալով թվային պարգև։ Ի տարբերություն վերահսկվող ուսուցման, ճիշտ պատասխանները նախապես տրված չեն. գործակալը պետք է ինքնուրույն գտնի այն գործողությունների հաջորդականությունը, որը առավելագույնի է հասցնում կուտակային պարգևը [աղբյուր՝ Sutton, Barto]։ Լուսացույցների կառավարումը այս ձևակերպման բնական օրինակ է. որոշումները (որ փուլը միացնել) ազդում են ապագա հերթերի վրա, և լավ որոշման արդյունքը երևում է ոչ թե անմիջապես, այլ որոշ ժամանակ անց։

\subsection{Մարկովյան որոշումների ընդունման պրոցես}

Ամրապնդող ուսուցման խնդիրը ձևակերպվում է որպես Մարկովյան որոշումների ընդունման պրոցես (Markov decision process, MDP), որը սահմանվում է հնգյակով $(S, A, P, r, \gamma)$։ Այստեղ $S$-ը վիճակների բազմությունն է, $A$-ն՝ գործողությունների բազմությունը, $P(s' \mid s, a)$-ն՝ անցումների հավանականությունը, $r(s,a)$-ն՝ պարգևի ֆունկցիան, իսկ $\gamma \in [0,1)$-ը՝ զեղչի գործակիցը։ Յուրաքանչյուր $t$ պահին գործակալը դիտարկում է վիճակը $s_t$, ընտրում է գործողություն $a_t$, ստանում է պարգև $r_{t+1}$ և անցնում նոր վիճակի $s_{t+1}$։

Գործակալի վարքագիծը նկարագրվում է քաղաքականությամբ $\pi(a \mid s)$՝ գործողությունների բաշխմամբ տվյալ վիճակում։ Նպատակը զեղչված կուտակային պարգևի մաթեմատիկական սպասումի առավելագույնացումն է.
\begin{equation}
  J(\pi) = \mathbb{E}_{\pi}\!\left[\sum_{k=0}^{\infty} \gamma^{k}\, r_{k+1}\right].
\end{equation}
Քաղաքականության որակը բնութագրվում է արժեքի ֆունկցիաներով.
\begin{equation}
  V^{\pi}(s) = \mathbb{E}_{\pi}\!\left[\sum_{k=0}^{\infty} \gamma^{k} r_{t+k+1} \,\Big|\, s_t = s\right],
  \qquad
  Q^{\pi}(s,a) = \mathbb{E}_{\pi}\!\left[\sum_{k=0}^{\infty} \gamma^{k} r_{t+k+1} \,\Big|\, s_t = s,\, a_t = a\right].
\end{equation}
Օպտիմալ գործողության արժեքի ֆունկցիան բավարարում է Բելմանի հավասարմանը.
\begin{equation}
  Q^{*}(s,a) = \mathbb{E}\!\left[\, r_{t+1} + \gamma \max_{a'} Q^{*}(s_{t+1}, a') \,\Big|\, s_t = s,\, a_t = a \right].
\end{equation}

\subsection{Q-ուսուցում}

Q-ուսուցումը (Q-learning) մեթոդ է, որը գնահատում է $Q^{*}$-ը՝ թարմացնելով աղյուսակի արժեքները փորձի հիման վրա.
\begin{equation}
  Q(s_t, a_t) \leftarrow Q(s_t, a_t) + \alpha\Bigl[\, r_{t+1} + \gamma \max_{a'} Q(s_{t+1}, a') - Q(s_t, a_t) \Bigr],
\end{equation}
որտեղ $\alpha$-ն ուսուցման արագությունն է [աղբյուր՝ Watkins]։ Մեթոդը աշխատում է ոչ միայն օպտիմալ, այլև ցանկացած քաղաքականության տվյալների հետ (off-policy)։ Նրա սահմանափակումն այն է, որ աղյուսակային ներկայացումը պահանջում է հնարավոր վիճակների ու գործողությունների թվարկում։ Լուսացույցների խնդրում վիճակը ներառում է հերթերի երկարությունները, որոնք անընդհատ մեծություններ են, ուստի աղյուսակային մեթոդները կիրառելու համար անհրաժեշտ է դիսկրետացում, ինչը կորցնում է ճշգրտությունը։

\subsection{Deep Q-Network (DQN)}

Այս սահմանափակումը հաղթահարվում է արժեքի ֆունկցիան նեյրոնային ցանցով մոտարկելով՝ $Q(s,a;\theta)$։ DQN մեթոդը [աղբյուր՝ Mnih et al., 2015] կայունացնում է ուսուցումը երկու մեխանիզմով։ Առաջինը փորձի վերարտադրման բուֆերն է (experience replay). անցումները $(s,a,r,s')$ պահվում են հիշողության մեջ և ուսուցման ժամանակ ընտրվում են պատահականորեն, ինչը նվազեցնում է հաջորդական տվյալների կոռելյացիան։ Երկրորդը թիրախային ցանցն է (target network)՝ $\theta^{-}$ պարամետրերով, որը թարմացվում է ավելի հազվադեպ։ Կորստի ֆունկցիան է.
\begin{equation}
  L(\theta) = \mathbb{E}\!\left[\Bigl(r + \gamma \max_{a'} Q(s', a'; \theta^{-}) - Q(s, a; \theta)\Bigr)^{2}\right].
\end{equation}
DQN-ը ենթադրում է դիսկրետ գործողությունների բազմություն, ինչը համապատասխանում է լուսացույցի փուլի ընտրությանը։

\subsection{Քաղաքականության գրադիենտ մեթոդներ և PPO}

Արժեքի վրա հիմնված մեթոդներին այլընտրանք են քաղաքականության գրադիենտ մեթոդները, որոնք օպտիմալացնում են պարամետրացված քաղաքականությունը $\pi_\theta(a \mid s)$ անմիջապես։ Գրադիենտը գնահատվում է որպես
\begin{equation}
  \nabla_\theta J(\theta) = \mathbb{E}\!\left[\nabla_\theta \log \pi_\theta(a \mid s)\, \hat{A}(s,a)\right],
\end{equation}
որտեղ $\hat{A}(s,a)$-ն առավելության ֆունկցիայի գնահատականն է. այն ցույց է տալիս, թե ընտրված գործողությունը որքանով է ավելի լավ կամ վատ վիճակի միջին արժեքից։ Գործնականում օգտագործվում է actor-critic կառուցվածքը՝ մեկ ցանցը (actor) որոշում է քաղաքականությունը, մյուսը (critic) գնահատում է վիճակի արժեքը։

Proximal Policy Optimization (PPO) մեթոդը [աղբյուր՝ Schulman et al., 2017] կայունացնում է ուսուցումը՝ սահմանափակելով քաղաքականության փոփոխությունը մեկ թարմացման ընթացքում։ Նշանակենք $\rho_t(\theta) = \pi_\theta(a_t \mid s_t)/\pi_{\theta_{\text{old}}}(a_t \mid s_t)$։ Այդ դեպքում նպատակային ֆունկցիան է.
\begin{equation}
  L^{\text{CLIP}}(\theta) = \mathbb{E}_t\!\left[\min\Bigl(\rho_t(\theta)\hat{A}_t,\; \operatorname{clip}\bigl(\rho_t(\theta),\, 1-\varepsilon,\, 1+\varepsilon\bigr)\hat{A}_t\Bigr)\right],
\end{equation}
որտեղ $\varepsilon$-ը կտրման պարամետրն է (գործնականում հաճախ $\varepsilon = 0{,}2$)։ PPO-ն համեմատաբար պարզ է իրականացնել, քիչ է զգայուն հիպերպարամետրերի նկատմամբ և լայնորեն կիրառվում է, այդ թվում՝ Stable-Baselines3 գրադարանում, որը օգտագործվում է սույն աշխատանքում։

\subsection{Եզրակացություն բաժնի վերաբերյալ}

Ամրապնդող ուսուցումը թույլ է տալիս կառավարման քաղաքականությունը ստանալ տվյալներից՝ առանց հոսքերի մաթեմատիկական մոդելի։ DQN-ը և PPO-ն հարմար են դիսկրետ գործողություններով խնդիրների համար։ Հաջորդ բաժնում դիտարկվում է, թե ինչպես են այս մեթոդները կիրառվել լուսացույցների կառավարման խնդրում։

\end{document}
